```latex
\documentclass[11pt]{article}

\usepackage{amsmath,amsthm,amssymb,enumitem,hyperref}

\newtheorem{theorem}{Theorem}
\newtheorem{lemma}[theorem]{Lemma}
\newtheorem{remark}[theorem]{Remark}

\begin{document}

\begin{center}
{\Large\bf Rigidity of a constant factor in a boundary family}
\end{center}

\medskip

We keep the notation of the problem.  Thus
\[
\partial V_t:=\overline{V_t}\cap \Delta_\Gamma\subset \Delta_\Gamma
\cong \mathcal M_{g_1,n_1}\times\cdots\times \mathcal M_{g_k,n_k},
\]
and for some fixed index $i$ we assume
\[
\pi_i(\partial V_t)=\{Y_t\}\qquad (t\in \Delta).
\]
We prove that $Y_t=Y_0$ for all $t$.

\medskip

As usual, one may pass to a sufficiently high finite level cover in order to avoid orbifold issues; this does not change the argument, and the conclusion descends.  We will not repeat this reduction below.

\begin{lemma}\label{lem:holomorphic-factor}
The map $t\mapsto Y_t$ is holomorphic.
\end{lemma}

\begin{proof}
By assumption \textup{(i)}, the sets $\partial V_t$ form a holomorphic family inside
$\Delta_\Gamma$; equivalently, after shrinking nothing, there is a complex submanifold
\[
\mathcal B\subset \Delta_\Gamma\times \Delta
\]
with fibers $\mathcal B_t=\partial V_t$, and with holomorphic projection
$\phi:\mathcal B\to \Delta$.

Consider the holomorphic map
\[
\Pi:=\pi_i\circ \mathrm{pr}_{\Delta_\Gamma}:\mathcal B\longrightarrow \mathcal M_{g_i,n_i}.
\]
By hypothesis, for each $t$, the restriction of $\Pi$ to the fiber
$\phi^{-1}(t)=\partial V_t$ is constant.  In local product coordinates for the smooth
holomorphic submersion $\phi$, this means precisely that $\Pi$ is independent of the
fiber variables.  Hence $\Pi$ factors through $\phi$, i.e. there exists a holomorphic
map
\[
Y:\Delta\longrightarrow \mathcal M_{g_i,n_i}
\]
such that $\Pi=Y\circ \phi$.  By construction, $Y(t)=Y_t$.
\end{proof}

\begin{lemma}\label{lem:teich-curve}
There exists an algebraic Teichm\"uller curve $T\subset V_0$ whose Deligne--Mumford
compactification $\overline{T}\subset \overline{\mathcal M}_g$ meets
$\partial V_0=\overline{V_0}\cap \Delta_\Gamma$ at a smooth point of $\partial V_0$.
\end{lemma}

\begin{proof}
This is a standard consequence of the structure theory of totally geodesic loci.

Indeed, by the theorem of Gekhtman--Markovi\'c / Eskin--Filip--Wright, the locus of
quadratic differentials whose Teichm\"uller disks are contained in the totally geodesic
surface $V_0$ is an affine invariant manifold $\mathcal N$.  The codimension-one
boundary piece of the Mirzakhani--Wright partial compactification of $\mathcal N$
lying over the curve $\partial V_0\subset \Delta_\Gamma$ is again an affine invariant
manifold; moreover, by the local product/collapse description of the boundary due to
Mirzakhani--Wright, a neighborhood of a smooth point of that boundary piece is obtained
by opening the collapsed cylinders.

Now closed $GL_2^+(\mathbb R)$-orbits are dense in affine invariant manifolds
(Eskin--Mirzakhani--Mohammadi).  Therefore one can choose a point of a closed
$GL_2^+(\mathbb R)$-orbit in $\mathcal N$ sufficiently near that boundary piece so that
collapsing the corresponding cylinders lands at a smooth point of the boundary piece
over $\partial V_0$.  The projection of that closed orbit to moduli space is an
algebraic Teichm\"uller curve $T\subset V_0$, and its stable limit under the cylinder
collapse is a point of $\partial V_0$.
\end{proof}

\begin{theorem}
Under the hypotheses of the problem, one has
\[
\pi_i(\partial V_t)=\{Y_0\}\qquad\text{for all }t\in \Delta.
\]
Equivalently, $Y_t=Y_0$ for every $t\in\Delta$.
\end{theorem}

\begin{proof}
By Lemma \ref{lem:holomorphic-factor}, it is enough to show that $Y_t$ is constant
on some neighborhood of $0$; the identity theorem will then imply that $Y$ is constant
on the connected disk $\Delta$.

Choose a Teichm\"uller curve $T\subset V_0$ as in Lemma \ref{lem:teich-curve}, and fix
a point
\[
p\in \overline{T}\cap \partial V_0\subset \Delta_\Gamma
\]
which is smooth on $\partial V_0$.

Since $\overline{\mathcal M}_g$ is projective and $\overline{T}$ is algebraic, Bertini
with prescribed base locus gives a hypersurface
\[
H\subset \overline{\mathcal M}_g
\]
such that:

\begin{enumerate}[label=\textup{(\alph*)}]
\item $\overline{T}\subset H$;
\item $H$ is transverse to $\overline{V_0}$ along $\overline{T}$;
\item after choosing a sufficiently small relatively compact neighborhood
$U\subset \overline{\mathcal M}_g$ of $\overline{T}$, one has
\[
H\cap U\cap \overline{V_0}=\overline{T}.
\]
\end{enumerate}

We now use assumptions \textup{(i)} and \textup{(ii)}.  Since $\overline{T}$ meets only
finitely many boundary strata, and on each stratum the intersections
$\overline{V_t}\cap \Delta_G$ form smooth families compatible under passage to deeper
strata, it follows (after shrinking the parameter disk to some smaller disk
$\Delta_\varepsilon\ni 0$) that the sets
\[
\overline{T}_t:=H\cap U\cap \overline{V_t}
\]
glue to a holomorphic family of compact curves in $\overline{\mathcal M}_g$, with
central fiber $\overline{T}_0=\overline{T}$.

Now we invoke M\"oller's rigidity theorem for Teichm\"uller curves
(\emph{Variations of Hodge structures of a Teichm\"uller curve}, JAMS 19 (2006)):
the stable map defined by the compactified Teichm\"uller curve $\overline{T}\to
\overline{\mathcal M}_g$ is rigid.  In particular, any holomorphic deformation of
$\overline{T}$ inside $\overline{\mathcal M}_g$ is locally trivial.  Applied to the
family $\{\overline{T}_t\}_{t\in \Delta_\varepsilon}$, this gives
\[
\overline{T}_t=\overline{T}\qquad\text{for all }t\in \Delta_\varepsilon.
\]

Hence the fixed point $p\in \overline{T}\cap \Delta_\Gamma$ belongs to
\[
\overline{T}_t\subset \overline{V_t}\cap \Delta_\Gamma=\partial V_t
\]
for every $t\in \Delta_\varepsilon$.  Since $\pi_i(\partial V_t)=\{Y_t\}$, we obtain
\[
Y_t=\pi_i(p)\qquad (t\in \Delta_\varepsilon).
\]
But $p\in \partial V_0$, and $\pi_i(\partial V_0)=\{Y_0\}$, so $\pi_i(p)=Y_0$.
Therefore
\[
Y_t=Y_0\qquad (t\in \Delta_\varepsilon).
\]

By Lemma \ref{lem:holomorphic-factor}, $Y:\Delta\to \mathcal M_{g_i,n_i}$ is holomorphic
and constant on the nonempty open set $\Delta_\varepsilon$.  Hence $Y$ is constant on
all of $\Delta$.  This is exactly the desired conclusion:
\[
\pi_i(\partial V_t)=\{Y_0\}\qquad\text{for all }t\in \Delta.
\]
\end{proof}

\bigskip

\noindent\textbf{References used in the proof.}
\begin{itemize}
\item M. M\"oller, \emph{Variations of Hodge structures of a Teichm\"uller curve},
J. Amer. Math. Soc. \textbf{19} (2006), 327--344.
\item A. Eskin, M. Mirzakhani, A. Mohammadi,
\emph{Isolation, equidistribution, and orbit closures for the $SL_2(\mathbb R)$ action
on moduli space}, Ann. of Math. \textbf{182} (2015), 673--721.
\item M. Mirzakhani, A. Wright,
\emph{The boundary of an affine invariant submanifold},
Invent. Math. \textbf{209} (2017), 927--984.
\item The identification of the cotangent locus of a totally geodesic subvariety with an
affine invariant manifold (work of Gekhtman--Markovi\'c and Eskin--Filip--Wright).
\end{itemize}

\end{document}
```